% Part: first-order-logic
% Chapter: natural-deduction
% Section: identity

\documentclass[../../../include/open-logic-section]{subfiles}

\begin{document}

\olfileid{fol}{ntd}{ide}

\olsection{له \usetoken{S}{identity} سره \usetoken{P}{derivation}}

له !!{identity} سره !!^{derivation}s اضافي استنتاجي قاعدو ته اړتيا لري۔

\begin{defish}
\AxiomC{}
\RightLabel{\Intro{\eq}}
\UnaryInfC{$\eq[t][t]$}
\DisplayProof
\hfill
\begin{tabular}{r}
\AxiomC{$\eq[t_1][t_2]$}
\AxiomC{$!A(t_1)$}
\RightLabel{\Elim{\eq}}
\BinaryInfC{$!A(t_2)$}
\DisplayProof
\\[3ex]
\AxiomC{$\eq[t_1][t_2]$}
\AxiomC{$!A(t_2)$}
\RightLabel{\Elim{\eq}}
\BinaryInfC{$!A(t_1)$}
\DisplayProof
\end{tabular}
\end{defish}

په پورته قاعدو کښې $t$، $t_1$ او~$t_2$ تړلې اصطلاحګانې دي۔
د~\Intro{\eq} قاعده راکوي چې د~$\eq[t][t]$ په بڼه هره د يووالي
جمله نېغ په نېغه، له هېڅ فرضيې پرته، !!{derive} کړو۔

\begin{ex}
که $s$ او~$t$ تړلې اصطلاحګانې وي، نو
$!A(s), \eq[s][t] \Proves !A(t)$:
\begin{prooftree}
\AxiomC{$\eq[s][t]$}
\AxiomC{$!A(s)$}
\RightLabel{$\Elim{\eq}$}
\BinaryInfC{$!A(t)$}
\end{prooftree}
دا به د «يو شان څيزونو د تعويض د اصل» يا د لايبنېڅ د قانون په نامه
اشنا وي۔
\end{ex}

\begin{prob}
ثابت کړئ چې $=$ متناظره او انتقالي دواړه ده؛ يعنې د
$\lforall[x][\lforall[y][(\eq[x][y] \lif
    \eq[y][x])]]$ او~$\lforall[x][\lforall[y][\lforall[z]((\eq[x][y]
    \land \eq[y][z]) \lif \eq[x][z])]]$
!!{derivation}s ورکړئ۔
\end{prob}

\begin{ex}
موږ دا !!{sentence}~!!{derive} کوو:
\begin{align*}
& \lforall[x][\lforall[y][((!A(x) \land !A(y)) \lif \eq[x][y])]]
\intertext{له دې !!{sentence} څخه:}
& \lexists[x][\lforall[y][(!A(y) \lif \eq[y][x])]]
\end{align*}
!!{derivation} له لاندې څخه پورته جوړوو:
\begin{prooftree}
\AxiomC{$\lexists[x][\lforall[y][(!A(y) \lif \eq[y][x])]]
\quad \Discharge{!A(a) \land !A(b)}{1}$}
\DeduceC{$\eq[a][b]$}
\DischargeRule{\Intro{\lif}}{1}
\UnaryInfC{$((!A(a) \land !A(b)) \lif \eq[a][b])$}
\RightLabel{\Intro{\lforall}}
\UnaryInfC{$\lforall[y][((!A(a) \land !A(y)) \lif \eq[a][y])]$}
\RightLabel{\Intro{\lforall}}
\UnaryInfC{$\lforall[x][\lforall[y][((!A(x) \land !A(y)) \lif \eq[x][y])]]$}
\end{prooftree}
اوس بايد اصلي فرضيه وکاروو: ځکه دا وجودي !!{formula} ده، د
منځنۍ پايلې~$\eq[a][b]$ د~!!{derive} کولو دپاره
\Elim{\lexists} کاروو۔
\begin{prooftree}
\AxiomC{$\lexists[x][\lforall[y][(!A(y) \lif \eq[y][x])]]$}
\AxiomC{$\Discharge{\lforall[y][(!A(y) \lif \eq[y][c]])}{2}$}
\noLine
\UnaryInfC{$\Discharge{!A(a) \land !A(b)}{1}$}
\DeduceC{$\eq[a][b]$}
\DischargeRule{\Elim{\lexists}}{2}
\BinaryInfC{$\eq[a][b]$}
\DischargeRule{\Intro{\lif}}{1}
\UnaryInfC{$((!A(a) \land !A(b)) \lif \eq[a][b])$}
\RightLabel{\Intro{\lforall}}
\UnaryInfC{$\lforall[y][((!A(a) \land !A(y)) \lif \eq[a][y])]$}
\RightLabel{\Intro{\lforall}}
\UnaryInfC{$\lforall[x][\lforall[y][((!A(x) \land !A(y)) \lif \eq[x][y])]]$}
\end{prooftree}
د پاس په ښي لوري کښې فرعي-!!{derivation} د خپلو فرضيو په کارولو
بشپړېږي، څو $\eq[a][c]$ او~$\eq[b][c]$ وښيي۔ دا دوه جلا
!!{derivation}s غواړي۔ د~$\eq[a][c]$~!!{derivation} داسې دے:
\begin{prooftree}
\AxiomC{$\Discharge{\lforall[y][(!A(y) \lif \eq[y][c]])}{2}$}
\RightLabel{\Elim{\lforall}}
\UnaryInfC{$!A(a) \lif \eq[a][c]$}
\AxiomC{$\Discharge{!A(a) \land !A(b)}{1}$}
\RightLabel{\Elim{\land}}
\UnaryInfC{$!A(a)$}
\RightLabel{\Elim{\lif}}
\BinaryInfC{$\eq[a][c]$}
\end{prooftree}
له~$\eq[a][c]$ او~$\eq[b][c]$ څخه
$\eq[a][b]$ د~\Elim{\eq} په وسيله !!{derive} کوو۔
\end{ex}

\begin{prob}
د لاندې !!{formula}s~!!{derivation}s ورکړئ:
\begin{enumerate}
\item $\lforall[x][\lforall[y][((\eq[x][y] \land !A(x)) \lif !A(y))]]$
\item $\lexists[x][!A(x)] \land \lforall[y][\lforall[z][((!A(y) \land
    !A(z)) \lif \eq[y][z])]] \lif \lexists[x][(!A(x) \land
  \lforall[y][(!A(y) \lif \eq[y][x])])]$
\end{enumerate}
\end{prob}

\end{document}
